% ═══════════════════════════════════════════════════════════════
% LERNBLATT: Reaktion von Metallen mit Sauerstoff
% Thema 2.4 | Chemie Klasse 8 | Kompakte 2-Seiten-Version
% ═══════════════════════════════════════════════════════════════

\documentclass[9pt, a4paper]{article}

\usepackage[utf8]{inputenc}
\usepackage[T1]{fontenc}
\usepackage[ngerman]{babel}

\usepackage{helvet}
\renewcommand{\familydefault}{\sfdefault}

\usepackage{microtype}
\usepackage[margin=1.2cm, top=1.4cm, bottom=1cm]{geometry}
\usepackage{enumitem}
\usepackage{tabularx}
\usepackage{booktabs}
\usepackage{array}
\usepackage{multicol}

\usepackage{amsmath, amssymb}
\usepackage{siunitx}
\sisetup{locale = DE, per-mode = symbol}

\usepackage{fancyhdr}
\usepackage{lastpage}

\usepackage{xcolor}
\usepackage{tcolorbox}
\tcbuselibrary{skins}

\usepackage{qrcode}

% ═══════════════════════════════════════════════════════════════
% FARBEN
% ═══════════════════════════════════════════════════════════════

\definecolor{chemie}{HTML}{2a7a4b}
\definecolor{hellgrau}{HTML}{F5F5F5}
\definecolor{mittelgrau}{HTML}{888888}
\definecolor{dunkelgrau}{HTML}{444444}

\colorlet{fachfarbe}{chemie}
\colorlet{fachfarbe-hell}{fachfarbe!15}

% ═══════════════════════════════════════════════════════════════
% KOPF- UND FUSSZEILE
% ═══════════════════════════════════════════════════════════════

\pagestyle{fancy}
\fancyhf{}
\renewcommand{\headrulewidth}{0pt}
\renewcommand{\footrulewidth}{0pt}
\fancyhead[L]{\small\textcolor{fachfarbe}{\textbf{Chemie}} | Thema 2.4}
\fancyhead[R]{\small Klasse 8}
\fancyfoot[C]{\small\textcolor{mittelgrau}{\thepage}}

% ═══════════════════════════════════════════════════════════════
% KOMPAKTE BOXEN
% ═══════════════════════════════════════════════════════════════

\newtcolorbox{formelkasten}{
    colback=fachfarbe-hell,
    colframe=fachfarbe,
    boxrule=1pt,
    arc=0pt,
    left=8pt, right=8pt, top=6pt, bottom=6pt,
    before skip=4pt, after skip=4pt
}

\newtcolorbox{merkkasten}{
    colback=hellgrau,
    colframe=hellgrau,
    boxrule=0pt,
    arc=0pt,
    left=6pt, right=6pt, top=5pt, bottom=4pt,
    borderline west={2pt}{0pt}{fachfarbe},
    before skip=4pt, after skip=4pt
}

% ═══════════════════════════════════════════════════════════════
% BEFEHLE
% ═══════════════════════════════════════════════════════════════

\newcommand{\seitentitel}[1]{%
    \noindent\textcolor{fachfarbe}{\rule{\textwidth}{1.5pt}}\\[-0.1em]
    {\Large\bfseries\textcolor{fachfarbe}{#1}}\\[-0.3em]
    \textcolor{fachfarbe}{\rule{\textwidth}{0.4pt}}
    \vspace{0.3em}
}

\newcommand{\abschnitt}[1]{%
    \vspace{6pt}%
    \noindent{\normalsize\bfseries\textcolor{fachfarbe}{#1}}%
    \vspace{2pt}%
    \nopagebreak[4]%
}

\setlength{\parindent}{0pt}
\setlength{\parskip}{2pt}
\setlength{\columnsep}{0.8cm}

% ═══════════════════════════════════════════════════════════════
% DOKUMENT
% ═══════════════════════════════════════════════════════════════

\begin{document}

% ═══════════════════════════════════════════════════════════════
% SEITE 1: GRUNDLAGEN + WERTIGKEITEN
% ═══════════════════════════════════════════════════════════════

\seitentitel{Reaktion von Metallen mit Sauerstoff}

\begin{multicols}{2}

\abschnitt{Was ist Oxidation?}

Die \textbf{Oxidation} ist die chemische Reaktion eines Stoffes mit Sauerstoff. Dabei entsteht ein \textbf{Oxid} und Energie wird frei (\textbf{exotherm}).

\vspace{4pt}
\begin{tabular}{@{}ll@{}}
\textbf{Element} & Stoff aus einer Atomsorte \\
\textbf{Verbindung} & Stoff aus verschiedenen Atomen \\
\textbf{Oxid} & Verbindung mit Sauerstoff \\
\end{tabular}

\abschnitt{Wortgleichungen}

\begin{formelkasten}
\centering
\textbf{Metall + Sauerstoff → Metalloxid}
\end{formelkasten}

\vspace{2pt}
\small
Mg + Sauerstoff → Magnesiumoxid\\
Fe + Sauerstoff → Eisenoxid\\
Al + Sauerstoff → Aluminiumoxid\\
Cu + Sauerstoff → Kupferoxid

\columnbreak

\abschnitt{Wertigkeiten}

Die \textbf{Wertigkeit} zeigt, wie viele Bindungen ein Atom eingeht.

\vspace{4pt}
\begin{center}
\small
\begin{tabular}{|c|c|c|c|c|c|c|}
\hline
\textbf{Mg} & \textbf{Al} & \textbf{Fe} & \textbf{Cu} & \textbf{Zn} & \textbf{Ca} & \textbf{O} \\
\hline
II & III & II/III & I/II & II & II & II \\
\hline
\end{tabular}
\end{center}

\abschnitt{Formelgleichung aufstellen}

\textbf{Beispiel:} Mg + O$_2$ → MgO

\textbf{1.} Formel: Mg(II) + O(II) → MgO (1:1)

\textbf{2.} Ausgleichen: Links 2 O, rechts 1 O!

\textbf{3.} Lösung: \textbf{2} Mg + O$_2$ → \textbf{2} MgO $\checkmark$

\end{multicols}

\begin{merkkasten}
\textbf{Merke:} Bei der Oxidation reagiert ein Element mit Sauerstoff zu einem Oxid. Die Reaktion ist exotherm – Energie (Wärme, Licht) wird frei.
\end{merkkasten}

\vspace{-2pt}
\hrule
\vspace{4pt}

% ═══════════════════════════════════════════════════════════════
% UNTERE HÄLFTE SEITE 1: BEOBACHTUNGEN + FLAMMENFARBEN
% ═══════════════════════════════════════════════════════════════

\begin{multicols}{2}

\abschnitt{Beobachtungen (Mg-Verbrennung)}

\small
\begin{tabular}{@{}p{1.8cm}p{5cm}@{}}
\textbf{Vorher} & Silbrig glänzendes Band \\
\textbf{Während} & Sehr helle, weiße Flamme \\
\textbf{Nachher} & Weißes Pulver (MgO) \\
\textbf{Energie} & Licht + Wärme (exotherm)
\end{tabular}

\columnbreak

\abschnitt{Flammenfarben der Metalle}

\small
\begin{tabular}{@{}lll@{}}
\textbf{Metall} & \textbf{Flamme} & \textbf{Oxid} \\
Mg & weiß & weiß \\
Fe & orange & braun \\
Al & silber-weiß & grau \\
Cu & grün-blau & schwarz
\end{tabular}

\end{multicols}

\abschnitt{Massenerhaltung}

\begin{formelkasten}
\centering
\textbf{Masse (Ausgangsstoffe) = Masse (Produkte)}
\end{formelkasten}

Das Oxid ist schwerer als das Metall, weil die Masse des Sauerstoffs aus der Luft hinzukommt.\\
\textbf{Beispiel:} 12\,g Mg + 8\,g O$_2$ → 20\,g MgO

% ═══════════════════════════════════════════════════════════════
% SEITE 2: PASSIVIERUNG + KORROSION
% ═══════════════════════════════════════════════════════════════

\newpage

\seitentitel{Reaktion von Metallen mit Sauerstoff}

\abschnitt{Oxidschichten: Porös oder dicht?}

\begin{center}
\small
\begin{tabular}{|p{3.5cm}|p{5cm}|p{5cm}|}
\hline
& \textbf{Magnesium} & \textbf{Aluminium} \\
\hline
\textbf{Oxidschicht} & \textbf{porös} (durchlässig) & \textbf{dicht} (undurchlässig) \\
\hline
\textbf{O$_2$ dringt durch?} & Ja → weitere Reaktion & Nein → Schutz \\
\hline
\textbf{Reaktion stoppt?} & Nein (vollständig) & Ja (nur Oberfläche) \\
\hline
\textbf{Fachbegriff} & Kettenreaktion & \textbf{Passivierung} \\
\hline
\end{tabular}
\end{center}

\begin{merkkasten}
\textbf{Passivierung:} Die dichte Oxidschicht schützt das Metall vor weiterer Korrosion. Deshalb „rostet" Aluminium nicht!
\end{merkkasten}

\abschnitt{Korrosion im Alltag}

\begin{tabular}{@{}p{2.5cm}p{12cm}@{}}
\textbf{Rost (Fe)} & Langsame Oxidation; poröse Schicht blättert ab → zerstört Metall \\
\textbf{Grünspan (Cu)} & Kupfercarbonat; schützt darunterliegendes Kupfer \\
\textbf{Al-Schicht} & Al$_2$O$_3$; dicht und fest → Passivierung
\end{tabular}

\vspace{8pt}

\begin{merkkasten}
\textbf{Das Wichtigste:}
\begin{itemize}[leftmargin=1.2em, itemsep=0pt, topsep=2pt]
    \item \textbf{Oxidation} = Reaktion mit O$_2$ → Oxid entsteht
    \item Reaktion ist \textbf{exotherm} (Energie frei)
    \item \textbf{Massenerhaltung:} Masse(Edukte) = Masse(Produkte)
    \item \textbf{Passivierung:} Dichte Schicht schützt (Al) vs. poröse Schicht (Mg, Fe)
\end{itemize}
\end{merkkasten}

% ═══════════════════════════════════════════════════════════════
% QR-CODES
% ═══════════════════════════════════════════════════════════════

\vfill

\begin{center}
\begin{tabular}{ccc}
\begin{minipage}{2.5cm}
\centering
\qrcode[height=1.5cm]{https://vimeo.com/699344409}\\
{\tiny Mg-Verbrennung}
\end{minipage}
&
\begin{minipage}{2.5cm}
\centering
\qrcode[height=1.5cm]{LP-01-metalloxide.html}\\
{\tiny Lernpfad}
\end{minipage}
&
\begin{minipage}{2.5cm}
\centering
\qrcode[height=1.5cm]{SIM-01-oxidation.html}\\
{\tiny Simulation}
\end{minipage}
\end{tabular}
\end{center}

\end{document}
